% =====================================================================
%  IEEE-STYLE CONFERENCE REPORT TEMPLATE  (self-contained)
%  MLE-AI-201 Machine Learning -- Final Project
%
%  This template MIMICS the IEEE two-column conference look using only
%  standard LaTeX packages. It does NOT depend on IEEEtran.cls, so it
%  compiles anywhere (tectonic, pdflatex, Overleaf) with nothing extra.
%
%  Build:   tectonic report.tex      (recommended)
%     or:   pdflatex report.tex ; pdflatex report.tex
% =====================================================================
\documentclass[10pt,twocolumn,a4paper]{article}

% ---- Geometry: IEEE-like margins / two narrow columns ----------------
\usepackage[a4paper,top=0.75in,bottom=1.0in,left=0.68in,right=0.68in,columnsep=0.25in]{geometry}

% ---- Times-like font (standard, ships with every TeX distro) ---------
\usepackage{mathptmx}          % Times text + math
\usepackage[T1]{fontenc}
\usepackage[utf8]{inputenc}    % allows accented / Azerbaijani chars

% ---- Common, dependency-free packages --------------------------------
\usepackage{graphicx}
\usepackage{amsmath,amssymb}
\usepackage{booktabs}
\usepackage{array}
\usepackage{titlesec}
\usepackage{caption}
\usepackage{cite}
\usepackage[hidelinks]{hyperref}
\usepackage{url}

% ---- IEEE-style section headings -------------------------------------
\renewcommand{\thesection}{\Roman{section}}
\renewcommand{\thesubsection}{\Alph{subsection}}
\titleformat{\section}[block]{\centering\normalsize\scshape}{\thesection.}{0.5em}{}
\titleformat{\subsection}[block]{\normalsize\itshape}{\thesubsection.}{0.5em}{}
\titlespacing*{\section}{0pt}{1.4ex plus 1ex minus .2ex}{1ex}
\titlespacing*{\subsection}{0pt}{1.0ex plus 1ex minus .2ex}{0.6ex}

% ---- Abstract / Index-Terms in the IEEE boldface-dash style ----------
\newcommand{\abstractkw}[1]{\noindent\textbf{\textit{#1}}}

\captionsetup{font=small,labelsep=period}
\setlength{\parindent}{1em}

\usepackage[safe]{tipa}
\DeclareUnicodeCharacter{0259}{\textschwa}
\DeclareUnicodeCharacter{018F}{\textschwa}
\input{numbers}

% =====================================================================
\begin{document}

\twocolumn[
  \begin{@twocolumnfalse}
  \begin{center}
    {\LARGE\bfseries Pricing Bina.az Listings from Scratch:\\Decision Trees and Pegasos SVMs vs.\ scikit-learn\par}
    \vspace{1.2em}
    {\large
      Allahverdi, Yunis\quad Gadirli, Ilaha\quad Kazimli, Fatime\par}
    \vspace{0.5em}
    {\itshape MLE-AI-201 Machine Learning, AI Academy --- Cohort II\par}
    \vspace{1.4em}
  \end{center}
  \begin{quote}
    \small
    \abstractkw{Abstract---}%
    We predict sale prices of Azerbaijani real-estate listings scraped from
    bina.az and frame the same data as two tasks: price regression
    (Task~A) and a premium/standard price-tier classification at the
    training median of \TierThreshold{}~AZN (Task~B). After deduplicating
    repeated scrapes and removing target-derived columns, \CleanRows{}
    listings remain. We implement a decision tree (Gini, entropy and MSE
    splits) and a soft-margin SVM trained with the Pegasos sub-gradient
    method, with a random-Fourier-feature extension, in NumPy only, and
    compare them with scikit-learn under matched settings. On the held-out
    test split our tree reaches $R^2=\TreeRtwo$ for price and F1~$=\TreeFone$
    for the tier, while the linear Pegasos SVM reaches F1~$=\SvmFone$
    (AUC~$\SvmAuc$). Our models track their scikit-learn counterparts
    closely, and tier error rates are highest for listings priced near the
    threshold.
    Code: \href{https://github.com/fatimekazimli/ML_Final_Projectcode}{\texttt{github.com/fatimekazimli/ML\_Final\_Projectcode}}
    (tag \texttt{v1.0-final}).

    \vspace{0.6em}
    \abstractkw{Index Terms---}decision tree, support-vector machine,
    Pegasos, random Fourier features, real-estate price prediction.
  \end{quote}
  \vspace{1.0em}
  \end{@twocolumnfalse}
]

% =====================================================================
\section{Introduction}
Listing prices on bina.az, Azerbaijan's largest real-estate platform, vary
by orders of magnitude with size, property type and location. We use a raw
scrape of these listings to study two classical models from the inside.
Our contributions are: (i) a documented, leakage-free cleaning and split of
a messy scraped dataset; (ii) a NumPy decision tree for both regression and
classification, verified against scikit-learn on toy data; (iii) a NumPy
Pegasos SVM with a random-Fourier-feature (RFF) extension, verified to
reach the same objective as \texttt{LinearSVC}; and (iv) a matched
benchmark against scikit-learn with cross-validated uncertainty and timing,
plus an error analysis. Every number in this report is produced by one
command, \texttt{python -m src.run\_all}.

\section{Background}
A decision tree recursively splits the feature space with axis-aligned
thresholds that maximise the decrease of an impurity measure (Gini or
entropy for classes, variance for regression). A soft-margin SVM finds the
hyperplane that minimises the regularised hinge loss; the Pegasos
algorithm~\cite{pegasos} solves this primal problem with stochastic
sub-gradient steps instead of a quadratic program. Random Fourier
features~\cite{rff} approximate an RBF kernel with an explicit
finite-dimensional map, so a linear solver can learn a non-linear boundary.

\section{Data and Preprocessing}
\label{sec:data}
% Data and Preprocessing -- Person 1 (Fatime Kazimli)
% every number here comes from report/numbers.tex, written by python -m src.run_all

We use version~1 of the Kaggle \emph{Binaaz sale project}
dataset\footnote{\url{https://www.kaggle.com/datasets/sehriyarmemmedli/binaaz-sale-project}},
a raw scrape of sale listings from bina.az (\texttt{house\_sale.csv},
downloaded 16~October~2026). It contains \RawRows{} rows. All cleaning is
implemented in \texttt{src/data\_prep.py}.

\paragraph{Cleaning}
The same listing is scraped repeatedly, so we keep only the latest scrape of
each listing URL (\UniqueRows{} unique listings). Prices are converted to AZN
at a fixed 1.70~AZN/USD rate (\UsdRows{} USD rows in this snapshot).
Text-encoded numbers are parsed: floor area (``85~m\textsuperscript{2}''),
rooms, floor and building height (``5~/~12'' gives \texttt{floor} and
\texttt{total\_floors}), and land area, where values in \emph{sot} are
converted to m\textsuperscript{2} (1~sot $=100$~m\textsuperscript{2}). We
apply fixed plausibility bounds chosen by hand, not estimated from data:
5{,}000--20{,}000{,}000~AZN for price and 10--2{,}000~m\textsuperscript{2}
for floor area when present. Rows with a missing price are dropped, never
imputed. \CleanRows{} listings remain.

\paragraph{Leakage}
Features are taken from an explicit allow-list (7 numeric: area, rooms,
floor, total floors, land area, latitude, longitude; 7 categorical:
building type, category, repair, mortgage, bill of sale, city, district),
so every other column is excluded by construction. The target-derived
columns we removed are: \LeakageCols{}. They are arithmetic functions of the
price and would make validation scores meaningless. Identifiers, URLs,
scrape timestamps and free text are also excluded. The full list of
excluded raw columns is: \ExcludedCols{}.

\paragraph{Encoding and split}
We split before computing any statistic. The split is stratified by fixed
AZN price bands (50k, 100k, 200k, 350k, 500k, 1M) with seed 42, giving
\TrainRows{} training, \ValRows{} validation and \TestRows{} test listings
(70/15/15). Numeric medians for imputation, the 25 most frequent districts
(the rest become ``other'') and the one-hot columns are learned on the
training rows only, giving \FeatureCount{} features. Features are
standardised with training mean and standard deviation for the SVM; trees
use them unscaled.

\paragraph{Two tasks}
Task~A predicts the sale price. Prices are heavy-tailed
(Fig.~\ref{fig:eda_price}), so models are trained on $\log(\text{price})$ and
scored with RMSE, MAE and $R^2$ in AZN (and on the log scale in
cross-validation). Task~B predicts a price tier: premium if the price is
above the training median of \TierThreshold{}~AZN, standard otherwise. The
threshold is fixed from the training split and applied unchanged to
validation and test, so \PremiumShareTrain\% of training listings are
premium (Fig.~\ref{fig:eda_balance}). Task~B is scored with precision,
recall, F1, ROC-AUC and the confusion matrix.

\paragraph{Exploration}
All exploratory plots use training rows only. Raw price has a long right
tail that the log transform largely removes (Fig.~\ref{fig:eda_price}),
which motivates modelling $\log(\text{price})$. Fig.~\ref{fig:eda_area}
shows floor area against price per property category, and
Fig.~\ref{fig:eda_map} shows prices on the map; floor area, longitude and
latitude are also among the most important tree features
(Sec.~\ref{sec:tree}), which is why the coordinates and district are kept.

\begin{table}[t]
  \centering
  \caption{Dataset summary after cleaning.}
  \label{tab:data}
  \begin{tabular}{@{}lr@{}}
    \toprule
    Property & Value \\
    \midrule
    Raw rows / unique listings  & \RawRows{} / \UniqueRows{} \\
    Rows after cleaning         & \CleanRows{} \\
    Features used               & \FeatureCount{} \\
    Train / Val / Test          & \TrainRows{} / \ValRows{} / \TestRows{} \\
    Tier threshold (AZN)        & \TierThreshold{} \\
    Premium share, train (\%)   & \PremiumShareTrain{} \\
    \bottomrule
  \end{tabular}
\end{table}

\begin{figure}[t]
  \centering
  \includegraphics[width=\linewidth]{figures/fig_price_hist.png}
  \caption{Training prices: raw AZN (left) and log scale with the tier threshold (right).}
  \label{fig:eda_price}
\end{figure}

\begin{figure}[t]
  \centering
  \includegraphics[width=0.85\linewidth]{figures/fig_area_vs_price.png}
  \caption{Floor area against price (log--log) for the five largest categories.}
  \label{fig:eda_area}
\end{figure}

\begin{figure}[t]
  \centering
  \includegraphics[width=0.85\linewidth]{figures/fig_map_latlng.png}
  \caption{Listing coordinates coloured by log(price).}
  \label{fig:eda_map}
\end{figure}

\begin{figure}[t]
  \centering
  \includegraphics[width=0.8\linewidth]{figures/fig_tier_balance.png}
  \caption{Standard vs premium counts in each split.}
  \label{fig:eda_balance}
\end{figure}


\section{Decision Tree}
\label{sec:tree}
\input{tree_section}

\section{Support-Vector Machine}
\label{sec:svm}
\input{svm_section}

\section{Results}
\label{sec:results}
Table~\ref{tab:results} reports held-out test metrics and
Table~\ref{tab:cv} the cross-validated mean $\pm$ std and timing for all
models. On Task~A our tree reaches RMSE \TreeRmse{}~AZN, MAE
\TreeMae{}~AZN and $R^2=\TreeRtwo$ on the test split. On Task~B our tree
reaches F1~$=\TreeFone$ (AUC~$\TreeAuc$), the linear Pegasos SVM
F1~$=\SvmFone$ (AUC~$\SvmAuc$) and the RFF SVM F1~$=\RffFone$
(AUC~$\RffAuc$). Each from-scratch model should be read against its
matched scikit-learn row rather than against the other model family. Our
implementations are slower to fit than the compiled scikit-learn code, as
the timing columns show.

\input{results_table}

\section{Discussion and Limitations}
Errors concentrate where the task is intrinsically ambiguous: for Task~B,
listings priced close to the tier threshold
(Sec.~\ref{sec:tree}, Fig.~\ref{fig:svm_errors}), and for Task~A, luxury
listings and land plots, whose AZN errors dominate RMSE. The trees are fit
on a seeded \TreeFitRows{}-row sample of the training split because the
Python implementation is slow, and the RBF \texttt{SVC} baseline is fit on
a 5{,}000-row subsample because its cost grows quadratically with $n$;
both choices are applied identically to our model and its matched
baseline. Impurity-based importances and SVM weights describe the fitted
models and are not causal effects. The data are one scrape of asking
prices, not transaction prices, and free-text descriptions were not used.
Next steps are ensembles of our trees, text features from the listing
descriptions, and a district-level model for sparse locations.

\section{Conclusion}
Implementing both models from first principles and checking them against
scikit-learn on toy problems and on real data showed that correct small
implementations reproduce library behaviour closely. Careful cleaning
(deduplicating scrapes, dropping target-derived columns, fixing every
statistic on the training split) mattered as much as the choice of model.

\section*{Contributions}
Allahverdi, Yunis: Pegasos SVM and RFF, SVM experiments, cross-validated
benchmark, slides. Gadirli, Ilaha: decision tree, tree experiments and
error analysis. Kazimli, Fatime: data cleaning and split, metrics, EDA,
\texttt{run\_all} integration and report assembly. Details are in
\texttt{contribution\_report.md}.

\section*{Tools and Acknowledgements}
AI assistants were used by the team: OpenAI Codex for the decision-tree
implementation, tests and its report section, and Anthropic Claude for the
SVM implementation, the data pipeline, integration code and drafting
report text. All members reviewed the AI-assisted code and text and are
responsible for it. Baselines use scikit-learn~\cite{sklearn}.

% =====================================================================
\begin{thebibliography}{9}
\bibitem{pegasos}
S.~Shalev-Shwartz, Y.~Singer, N.~Srebro, and A.~Cotter,
``Pegasos: primal estimated sub-gradient solver for SVM,''
\emph{Mathematical Programming}, vol.~127, no.~1, pp.~3--30, 2011.
\bibitem{rff}
A.~Rahimi and B.~Recht, ``Random features for large-scale kernel machines,''
in \emph{Advances in Neural Information Processing Systems}, vol.~20, 2007.
\bibitem{sklearn}
F.~Pedregosa \emph{et al.}, ``Scikit-learn: Machine learning in Python,''
\emph{Journal of Machine Learning Research}, vol.~12, pp.~2825--2830, 2011.
\end{thebibliography}

\end{document}
